\section{Discussion}
\label{sec:discussion}

\subsection{Limitations}
\label{subsec:limitations}

\paragraph{Confident wrongness.} \eps cannot detect a model that is
wrong and confident. If training placed all mass on a deprecated API,
the token probability is sharp and \eps is low; the system emits the
wrong code with \textsc{complete} status. Fundamental limit of any
generation-time signal; motivates static-analysis integration.

\paragraph{Confidently-resolved domains.} Prompts where every model
has converged on a strong training-data prior produce $\eps = 0.000$
regardless of correctness. Current-UTC datetime retrieval and
HTTP-GET-to-JSON fall here across all four models. For stable
widely-deployed APIs this is desirable; for deprecated-but-%
confidently-recalled APIs it is Scenario B's failure mode.

\paragraph{Small-sample calibration.} Main calibration is 30 prompts
per model; Wilson intervals on HIGH detection are $\pm 17$pp at
$n=20$. We claim recall-first design plus an empirical miss rate of
zero on the sets evaluated, not a proof of perfect recall.

\paragraph{Type 3 floor.} Implementation-approach uncertainty on
simple prompts (brute vs.\ hash two-sum) legitimately raises \eps on
correct code. AST filtering does not remove it; the $\sim 50\%$ LOW
FP is predominantly Type 3 and unlikely to shrink without sacrificing
recall.

\paragraph{LLM-judged ground truth.}\label{subsec:gt-disclosure}
LOW ground truth is \texttt{exec()}-verified against assertions covering
return type, edge cases, and correctness. HIGH ground truth is LLM-judged
via GPT-4o. Circularity exists where the evaluated model shares a family
with the judge (GPT-4o-mini judged by GPT-4o); results for that model
should be read with that caveat. We disclose rather than overclaim human
verification we did not perform.

\paragraph{Hyperparameters.} The cascaded tier algorithm (six checks,
three thresholds) was calibrated on GPT-4o; other models inherit
these. The binary \textsc{flagged+} gate is robust; internal tier
placement is a refinement on top.

\subsection{Context learning as a microcosm}
\label{subsec:context}

Injecting learned false-positive patterns into the reviewer's context
(the configuration evaluated in \S\ref{subsec:review-loop}) raised
clearance from $53.2\%$ to $68.7\%$, and late false positives fell from
$4$ to $0$. This is in-context improvement on a fixed model;
training-time improvement would be expected to compound with it.

\subsection{Future work}
\label{subsec:future}

\textbf{Pipeline propagation:} upstream \eps should propagate as
prior uncertainty into downstream generations, modulating sampling
temperature or review threshold. \textbf{Open-source models:} vLLM
exposes full top-$k$; a Llama/Qwen/Mistral calibration run would test
whether the recall-first boundary holds outside OpenAI. \textbf{Per-%
token intervention:} surface the high-\eps token as an inline editor
marker, or trigger targeted re-decoding with lower temperature on
just the uncertain span. \textbf{Cross-decision consistency:}
Scenario D's mismatches argue for AST-level consistency checks over
\eps (e.g., \texttt{async def} $+$ blocking HTTP client).
\textbf{Import-gated activation:} restricting \eps scoring to
functions that import external libraries (i.e., skipping pure-logic
functions) would eliminate the Type~3 false-positive floor at
negligible cost, since production API-integration code almost always
begins with an import.
